L'objectif de ce projet est de concevoir un orchestrateur cloud intelligent pour une Smart City, capable de trouver le compromis parfait entre les temps de réponse garantis (SLA) et la réduction de l'empreinte énergétique des serveurs.
Pour y parvenir, nous avons développé un agent d'apprentissage par renforcement (PPO) qui a appris à piloter simultanément la quantité de services allumés (Autoscaling) et leur placement sur le matériel (Scheduling) en surmontant les pièges du gaspillage et du démarrage à froid.
Ce travail s'intègre au sein d'un \textit{benchmark} global confrontant cette approche PPO à une stratégie classique (Hybride), un modèle par valeur (DQN) et une architecture décentralisée (Multi-Agent) afin d'identifier le meilleur système de décision pour l'avenir.